% Generated draft: requires mathematical and source review.
\documentclass[11pt]{article}
\usepackage[margin=25mm]{geometry}
\usepackage{fontspec}
\setmainfont{Latin Modern Roman}
\usepackage{amsmath,amssymb,mathtools,unicode-math}
\setmathfont{Latin Modern Math}
\usepackage{enumitem,needspace,xurl,hyperref}
\setlength{\parindent}{0pt}
\setlength{\parskip}{5pt}
\setlength{\emergencystretch}{3em}
\newcommand{\N}{\mathbb{N}}
\newcommand{\Z}{\mathbb{Z}}
\newcommand{\Q}{\mathbb{Q}}
\newcommand{\R}{\mathbb{R}}
\newcommand{\C}{\mathbb{C}}
\newcommand{\F}{\mathbb{F}}
\newcommand{\A}{\mathbb{A}}
\newcommand{\PP}{\mathbb{P}}
\newcommand{\E}{\mathbb{E}}
\newcommand{\eps}{\varepsilon}
\newcommand{\dd}{\,\mathrm{d}}
\begin{document}
\section*{UnsolvedMath ID: 8800054}
\textbf{Source problem code:} catalogue entry
\subsection{Definitions and mathematical statement}
A hyperelliptic curve \(C/\F_q\) is, in the usual genus-at-least-two convention, a smooth projective geometrically connected curve admitting a degree two map to \(\mathbb P^1\). Its Jacobian \(J_C\) is the abelian variety of degree-zero divisor classes. For a prime \(r\mid\#J_C(\F_q)\) with \(r\nmid q\), the embedding degree is
\[
 k=\operatorname{ord}_r(q)=\min\{j\geq1:r\mid q^j-1\}.
\]
Then the \(r\)-th roots of unity lie in \(\F_{q^k}^{\times}\). Pairing constructions use suitable torsion subgroups or quotient groups and a nondegenerate bilinear map to these roots of unity; the two pairing inputs need not be the same rational cyclic subgroup.

Lange asks to construct pairing-friendly hyperelliptic curves suitable for cryptography. In the customary interpretation, the construction should provide curves and explicit subgroup parameters with a large prime \(r\) dividing the Jacobian order, a manageable embedding degree, and efficient group and pairing operations at the intended security level. These are qualitative design goals, not numerical constraints given in the excerpt.

The source specifies no genus, characteristic, field-size range, subgroup-size ratio, embedding degrees, security target, or efficiency benchmark. It also does not state whether individual examples or infinite parameterized families are required. Thus no claim of construction for every prescribed embedding degree or of efficiency matching elliptic curves is added. The supplied summary reports partial genus-two constructions, not a complete answer to every such possible design specification.
\subsection{Short English statement}
Construct hyperelliptic curves whose Jacobians offer large prime-order subgroups and practical cryptographic pairings. The source leaves genus, embedding degrees, field sizes, efficiency, and security criteria qualitative.
\subsection{Sources}
\begin{itemize}
\item Local UnsolvedMath registry, permanent ID 8800054, source code catalogue entry; attributed to Tanja Lange; no proposal year supplied; update 2026-08-06.
\item IPAM, \emph{Workshop on Number Theory and Cryptography --- Open Problems} (2006), attachment 6239, PDF page 41, printed slide 39: \url{https://www.ipam.ucla.edu/programs/workshops/workshop-i-number-theory-and-cryptography-open-problems/?tab=schedule}. Registry access date 2026-07-29; PDF and legacy-PPT visual audit. Provenance reports 19 attachments, 17 PDFs and two legacy PPT decks, with repetitions and solved definitions excluded. Present status and redistribution terms require review.
\item The supplied partial-progress summary mentions Kawazoe--Takahashi, Freeman, and genus-two complex-multiplication constructions without exact references or parameter/efficiency theorems. Its comparison with elliptic pairing families is an unverified registry assessment.
\end{itemize}
\end{document}
